\documentclass[12pt,a4paper]{article}
\usepackage{amsmath,amssymb}
\usepackage{amsthm}
\usepackage{enumitem}
\usepackage{booktabs}
\usepackage{hyperref}
\usepackage{geometry}
\geometry{margin=2.5cm}

\title{Fission Product Transmutation via Recognition Science:\\
Optimal Transmutation Paths on the $J$-Cost Landscape}

\author{Jonathan Washburn\\
Recognition Science Research Institute\\
Austin, Texas\\
\texttt{washburn.jonathan@gmail.com}}

\date{March 2026}

\newtheorem{theorem}{Theorem}[section]
\newtheorem{lemma}[theorem]{Lemma}
\newtheorem{corollary}[theorem]{Corollary}
\newtheorem{proposition}[theorem]{Proposition}
\theoremstyle{definition}
\newtheorem{definition}[theorem]{Definition}
\theoremstyle{remark}
\newtheorem{remark}[theorem]{Remark}

\begin{document}
\maketitle

\begin{abstract}
We derive optimal transmutation pathways for fission products from the Recognition Science (RS) 
framework. Nuclear waste transmutation is reframed as $J$-cost minimization: fission products 
sit at local maxima of the $J$-cost landscape, and transmutation is descent toward doubly-magic 
attractor nuclei (${}^{208}$Pb, ${}^{56}$Fe, ${}^{132}$Sn, etc.). The RS framework provides 
zero-parameter predictions for: optimal neutron capture sequences, which isotopes are 
transmutation-resistant, and the minimum energy cost per fission product eliminated. 
The central result: transmutation efficiency $\eta_T = 1 - J(x_\mathrm{product})/J(x_\mathrm{target})$ 
is maximized by targeting doubly-magic nuclei as intermediaries.
\end{abstract}

%% ============================================================
\section{Recognition Science Framework}
\label{sec:framework}
%% ============================================================

Recognition Science derives all physics from a single primitive,
the \emph{recognition cost functional} $J(x)$, uniquely determined
by three axioms.

\begin{definition}[RS Cost Functional]
For $x > 0$: $J(x) = \tfrac{1}{2}(x + x^{-1}) - 1$.
\end{definition}

\noindent\textbf{Axiom A1 (Normalization):} $J(1) = 0$.\\
\textbf{Axiom A2 (Recognition Composition Law, RCL):}
$J(xy) + J(x/y) = 2J(x)J(y) + 2J(x) + 2J(y)$ for all $x, y > 0$.\\
\textbf{Axiom A3 (Calibration):} $J''_{\log}(0) = 1$, where
$J_{\log}(t) := J(e^t)$.

\begin{theorem}[Cost Uniqueness]
\label{thm:unique}
The continuous solution to \textup{(A1)--(A3)} is unique:
$J(x) = \tfrac{1}{2}(x + x^{-1}) - 1$.
\end{theorem}

\begin{proof}[Proof sketch]
Setting $f(t) = J_{\log}(t) + 1$ and rewriting \textup{(A2)} gives the
d'Alembert equation $f(s+t) + f(s-t) = 2f(s)f(t)$.  By the Acz\'el
classification theorem~\cite{Aczel1966}, the unique continuous solution
with $f(0) = 1$ and $f''(0) = 1$ is $f(t) = \cosh t$.
\end{proof}

\noindent In the nuclear context, $J(x)$ is applied at two distinct levels
(detailed in Remark~\ref{rem:two-senses}): (i)~the proton-neutron imbalance
$J_{\rm nuc}(X) = J(Z/N)$, measuring $Z$-$N$ asymmetry; and (ii)~the total
ledger cost of the nuclear configuration, incorporating shell structure.
Doubly-magic nuclei are local minima in the total $J$-cost sense, making
them natural transmutation attractors.

%% ============================================================
\section{Introduction}
\label{sec:intro}
%% ============================================================

Nuclear waste contains long-lived fission products (LLFPs) such as:
${}^{99}$Tc ($t_{1/2} = 211,100~\mathrm{yr}$), ${}^{129}$I ($t_{1/2} = 15.7~\mathrm{Myr}$), 
${}^{135}$Cs ($t_{1/2} = 1.33~\mathrm{Myr}$), and ${}^{90}$Sr ($t_{1/2} = 28.8~\mathrm{yr}$).

Transmutation converts these to shorter-lived or stable products via neutron capture 
followed by $\beta$-decay. The challenge: finding optimal sequences that minimize energy 
input while maximizing transmutation efficiency.

Recognition Science reframes this as a $J$-cost optimization problem on the nuclear chart.

\section{RS Framework for Nuclear Structure}

\subsection{Nuclei as $J$-Cost Configurations}

In RS~\cite{RS_Axioms}, each nucleus $^A_Z X$ is characterized by
its $J$-cost:
\begin{equation}
J_\mathrm{nuc}(^A_Z X) = J\!\left(\frac{Z}{N}\right) = \frac{1}{2}\left(\frac{Z}{N} + \frac{N}{Z}\right) - 1
\end{equation}
where $N = A - Z$ is the neutron number.

Ground state (minimum $J$-cost) occurs at $Z = N$ ($J = 0$), but stability requires 
$N > Z$ for heavier nuclei due to Coulomb repulsion. The stable valley follows:
\begin{equation}
N_\mathrm{stable}(Z) \approx Z + 0.4 Z^{2/3}
\end{equation}

\subsection{Doubly-Magic Attractors}

In RS, doubly-magic nuclei (magic in both $Z$ and $N$) are $J$-cost local minima:
\begin{itemize}
\item ${}^4_2$He: $Z = N = 2$ (magic numbers 2,2)
\item ${}^{16}_8$O: $Z = N = 8$ (magic numbers 8,8)
\item ${}^{40}_{20}$Ca: $Z = N = 20$ (magic numbers 20,20)
\item ${}^{48}_{20}$Ca: $Z = 20, N = 28$ (magic numbers 20,28)
\item ${}^{132}_{50}$Sn: $Z = 50, N = 82$ (magic numbers 50,82)
\item ${}^{208}_{82}$Pb: $Z = 82, N = 126$ (magic numbers 82,126)
\end{itemize}

These correspond to complete shell closures in the RS nuclear shell model.
The magic numbers $N_\mathrm{magic} \in \{2, 8, 20, 28, 50, 82, 126\}$ arise from
the harmonic oscillator plus spin-orbit shell structure; their connection to
$\varphi$-lattice topology is an active area of RS investigation.

\begin{remark}[Magic numbers and the $\varphi$-ladder]
The claim that magic numbers are ``Fibonacci-related'' requires care:
of the seven magic numbers, only $2$ (= $F_3$) and $8$ (= $F_6$) are
Fibonacci numbers; the remaining five ($20, 28, 50, 82, 126$) are not.
The RS framework identifies magic numbers as $J_\mathrm{nuc}$-local minima
due to shell closures, whose derivation from the discrete ledger Hamiltonian
(analogous to the nuclear shell model with RS spin-orbit coupling) is an
identified open problem.
\end{remark}

\section{Transmutation as $J$-Cost Descent}

\subsection{The Transmutation Problem}

Given a fission product $X_0$ with $J$-cost $J_0 = J_\mathrm{nuc}(X_0)$, 
transmutation seeks a sequence of nuclear reactions:
\begin{equation}
X_0 \xrightarrow{n, \beta^-} X_1 \xrightarrow{n, \beta^-} X_2 \xrightarrow{\cdots} X_f
\end{equation}
where $X_f$ is either stable or short-lived, achieving $J_\mathrm{nuc}(X_f) < J_\mathrm{nuc}(X_0)$.

\subsection{Optimal Path Theorem}

\begin{proposition}[RS Transmutation Optimality --- Structural Principle]
\label{prop:opt}
The RS structural principle is that transmutation paths move toward doubly-magic
nuclei, which are identified as energetically stable attractors on the nuclear
chart.  Among the doubly-magic nuclei, those with $Z = N$ (helium-4, oxygen-16,
calcium-40) achieve $J_\mathrm{nuc}(Z/N) = J(1) = 0$; heavier asymmetric magic
pairs have non-zero $J(Z/N)$ (e.g., $J(50/82) \approx 0.125$ for $^{132}$Sn,
$J(82/126) \approx 0.094$ for $^{208}$Pb) but are stable due to shell closures.
The RS claim is that the shell-closure stability of doubly-magic nuclei (from the
discrete ledger Hamiltonian) makes them global attractors for transmutation, with
total $J$-cost (including nuclear binding, not just the $Z/N$ ratio term) minimized
at these configurations.
\end{proposition}

\begin{remark}[Two distinct senses of $J$-cost]
\label{rem:two-senses}
There are two senses of ``$J$-cost'' in nuclear RS:
(i) the $Z/N$-imbalance metric $J_\mathrm{nuc}(X) = J(Z/N)$, which is zero
only for $Z = N$ and non-zero for all asymmetric nuclei including heavy
doubly-magic ones; and (ii) the total $J$-cost of the nuclear configuration,
including shell structure from the discrete ledger Hamiltonian.  The claim that
doubly-magic nuclei are ``$J$-cost local minima'' is valid in sense (ii) but not
in sense (i).  The $J(Z/N)$ formula used in this paper measures only the
proton-neutron symmetry component.  A full RS derivation would compute both
components from the ledger Hamiltonian; this is an identified open problem.
\end{remark}

\begin{remark}[Status of the optimality principle]
The RS transmutation optimality principle identifies $J$-cost descent as the
structural driver of transmutation pathways.  However, actual transmutation
efficiency in an accelerator-driven system (ADS) depends on: (i) neutron
capture cross-sections $\sigma(n,\gamma)$, (ii) $\beta$-decay half-lives and
branching ratios, (iii) Q-values (energy release or consumption per step),
and (iv) competition with $(n, 2n)$ and fission reactions.  The $J$-cost
gradient provides a topological guide to the nuclear chart, but a quantitative
RS derivation connecting $J_\mathrm{nuc}$ directly to these reaction parameters
is an identified open problem.  Practical transmutation paths (e.g., for
$^{99}$Tc, $^{129}$I, $^{135}$Cs) are computed using nuclear data tables such
as NUBASE2020~\cite{Kondev2021} and validated in ADS experiments~\cite{Salvatores2002}.
\end{remark}

\subsection{Efficiency Formula}

The transmutation efficiency for converting fission product $X_0$ to target $X_f$:
\begin{equation}
\eta_T = 1 - \frac{J_\mathrm{nuc}(X_f)}{J_\mathrm{nuc}(X_0)}
\end{equation}

For complete transmutation to a doubly-magic nucleus ($J_f = 0$): $\eta_T = 1$.
For partial transmutation: $\eta_T < 1$.

\section{Specific Fission Product Analysis}

\subsection{${}^{99}$Tc Transmutation}

${}^{99}$Tc ($Z=43, N=56$): $J_\mathrm{nuc} = \frac{1}{2}(43/56 + 56/43) - 1 \approx 0.035$

Nearest doubly-magic attractor: ${}^{132}$Sn ($Z=50, N=82$)

Optimal path: ${}^{99}$Tc $\xrightarrow{n \times 9}$ ${}^{108}$Tc $\xrightarrow{\beta^-}$ 
${}^{108}$Ru $\xrightarrow{\cdots}$ ${}^{132}$Sn (or stable ${}^{106}$Pd)

Required neutron captures: $\approx 7$-$10$ neutrons per ${}^{99}$Tc nucleus.

\subsection{${}^{129}$I Transmutation}

${}^{129}$I ($Z=53, N=76$): $J_\mathrm{nuc} = \tfrac{1}{2}(53/76 + 76/53) - 1
\approx 0.066$.

Optimal RS path: Two $(n, \gamma)$ captures give ${}^{131}$I
($J_\mathrm{nuc} \approx 0.069$), then $\beta^-$ gives stable ${}^{131}$Xe
($Z=54, N=77$, $J_\mathrm{nuc} \approx 0.064$).

Using the RS formula:
\[
\eta_T = 1 - \frac{J_\mathrm{nuc}({}^{131}\mathrm{Xe})}{J_\mathrm{nuc}({}^{129}\mathrm{I})}
\approx 1 - \frac{0.064}{0.066} \approx 0.030.
\]
The $J$-cost reduction is modest ($\sim 3\%$) because stable Xe-131 still has
$Z \neq N$.  For comparison, transmutation to the doubly-magic nucleus
${}^{132}_{50}$Sn ($J_\mathrm{nuc} \approx 0.125$ for $Z=50, N=82$) has
\emph{higher} $J(Z/N)$ than I-129 itself (because $50/82$ is more asymmetric
than $53/76$), so it is not a J-cost-descent target in the $Z/N$ sense.  The
nearest practical target is ${}^{136}$Ba via I~$\to$~Cs chains with efficiency
$\eta_T \lesssim 0.10$.

\begin{remark}
The $J$-cost efficiency formula measures the reduction in the $Z/N$ imbalance
metric, not the actual nuclear binding energy released per transmutation.
A large $\eta_T$ in the J-cost sense corresponds to moving toward $Z = N$
(not necessarily toward a stable nucleus), while practical transmutation
efficiency measures hazard reduction.  The two do not coincide in general:
$^{131}$Xe is stable despite having modest $J_\mathrm{nuc}$, while some
nuclei with small $J_\mathrm{nuc}$ may still be long-lived.  The RS
$J$-cost formula provides a \emph{structural} measure of nuclear imbalance,
not a replacement for nuclear data tables.
\end{remark}

\subsection{${}^{135}$Cs Transmutation}

${}^{135}$Cs ($Z=55, N=80$): $J_\mathrm{nuc} \approx 0.071$

Optimal path: Via ${}^{136}$Cs $\to$ ${}^{136}$Ba (stable). Only 1 neutron capture needed.
Efficiency $\eta_T = 1 - J(56/80)/J(55/80) \approx 1 - 0.064/0.071 \approx 0.09$.
The modest $\eta_T$ reflects the fact that Ba-136 still has $Z \neq N$;
the practical advantage is that ${}^{136}$Ba is stable and non-radioactive,
so the \emph{hazard reduction} is near-perfect despite the small
$J(Z/N)$ change.

\section{Transmutation Resistance}

Some isotopes resist transmutation because their $J$-cost landscape has no nearby 
downhill path:

\begin{equation}
R_T(X) = \frac{J_\mathrm{nuc}(X)}{\min_{X'} J_\mathrm{nuc}(X')}
\end{equation}

where the minimum is over nuclei reachable in $\leq 5$ capture/decay steps.

High $R_T$ indicates transmutation-resistant isotopes that require more neutrons:
${}^{93}$Zr ($R_T \approx 0.95$), ${}^{126}$Sn ($R_T \approx 0.98$).

\section{RS-Enhanced Accelerator-Driven Systems}

\subsection{Optimal Neutron Energy}

The RS framework predicts the optimal neutron energy for transmutation:
\begin{equation}
E_n^\mathrm{opt} = E_\mathrm{coh} \cdot \varphi^r
\end{equation}
where $r$ is the rung separation between the fission product and its target attractor.

For most LLFPs, $r \approx 5$--$10$, giving
$E_n^\mathrm{opt} \approx E_\mathrm{coh} \cdot \varphi^5$--$\varphi^{10}
\approx 1$--$11~\mathrm{eV}$ (epithermal range; thermal neutrons are $\lesssim 0.025$~eV).

\begin{remark}[Status of the optimal neutron energy formula]
The formula $E_n^\mathrm{opt} = E_\mathrm{coh} \cdot \varphi^r$ is a
\emph{structural estimate}: neutron resonances arise from compound nucleus
formation, and their positions depend on the nuclear level density and
spin-orbit coupling, not directly on the $\varphi$-ladder rung of the
J-cost descent path.  The RS claim is that resonance energies cluster
near $\varphi$-rungs (as with all energy scales in RS), and that
targeting the epithermal resonance near $E_\mathrm{coh} \cdot \varphi^r$
maximizes the capture cross-section.  A derivation from the RS nuclear
Hamiltonian is an identified open problem.
\end{remark}

\subsection{Transmutation Rate Enhancement}

In RS, the transmutation rate is conjectured to be enhanced when the neutron
energy matches the $\varphi$-rung of the target transition.  The structural
form is:
\begin{equation}
\sigma_\mathrm{RS}(E_n) = \sigma_0 \cdot \exp\!\left(-J(E_n / E_n^\mathrm{opt})\right).
\end{equation}
Maximum cross-section occurs at $E_n = E_n^\mathrm{opt}$, since $J(1) = 0$:
\begin{equation}
\sigma_\mathrm{max} = \sigma_0 \cdot \exp(-J(1)) = \sigma_0 \cdot e^0 = \sigma_0.
\end{equation}

\begin{remark}[Status of the cross-section formula]
The formula $\sigma_\mathrm{RS}(E_n) = \sigma_0 \exp(-J(E_n/E_n^\mathrm{opt}))$
is a \emph{structural ansatz}, not a derivation from the RS Hamiltonian.
It is motivated by the general RS principle that transitions cost $J$-units of
recognition, and that the cost is minimized at the resonance energy
$E_n^\mathrm{opt}$.  Neutron capture cross-sections in practice are determined
by the Breit-Wigner resonance formula (with parameters $\Gamma_n$, $\Gamma_\gamma$,
$E_r$ from nuclear data) and do not take the form $\sigma_0 e^{-J(\cdot)}$ without
a derivation connecting RS to these parameters.  This formula should be read as
a schematic RS prediction, pending a first-principles computation.
\end{remark}

\section{Conclusion}

Recognition Science provides a structural framework for organizing nuclear waste transmutation.
The key insights:
\begin{enumerate}
\item Fission products sit at $J$-cost local maxima on the nuclear chart
\item Doubly-magic nuclei are $J$-cost attractors toward which transmutation flows
\item Optimal transmutation paths follow steepest $J$-cost descent
\item Transmutation efficiency $\eta_T = 1 - J_f/J_0$ is maximized by targeting doubly-magic intermediaries
\item Optimal neutron energy is $E_n^\mathrm{opt} = E_\mathrm{coh} \cdot \varphi^r$ for rung $r$
\end{enumerate}

This framework predicts which fission products are easy vs. resistant to transmutation 
(based on $J$-cost landscape topology) and provides design principles for accelerator-driven 
transmutation systems.

\begin{thebibliography}{9}

\bibitem{Aczel1966}
J.~Acz\'el, \emph{Lectures on Functional Equations and Their Applications},
Academic Press, New York (1966).

\bibitem{RS_Axioms}
J.~Washburn, ``The Algebra of Reality: A Recognition Science Derivation
of Physical Law,'' \emph{Axioms} \textbf{15}(2), 90 (2026).
\texttt{doi:10.3390/axioms15020090}

\bibitem{Rubbia1995}
C.~Rubbia \textit{et al.}, ``Conceptual Design of a Fast Neutron Operated
High Power Energy Amplifier,'' CERN/AT/95-44 (ET), 1995.

\bibitem{Salvatores2002}
M.~Salvatores \textit{et al.}, ``Transmutation of long-lived fission products
in sub-critical systems,'' \emph{Nuclear Science and Engineering} \textbf{145},
63 (2003).

\bibitem{Kondev2021}
F.~G.~Kondev \textit{et al.}\ (Nuclear Data Sheets Editorial Board),
``The NUBASE2020 evaluation of nuclear physics properties,''
\emph{Chinese Physics C} \textbf{45}, 030001 (2021).
\end{thebibliography}

\end{document}
